\section{Derivations}

\subsection{Markov equilibrium}\label{app:markov}

From the market clearing for bonds and the expression for $B_{i,t}$ from the Lemma \ref{lemma: investor_problem}, we obtain
\begin{equation}\nonumber
    \sum_{i=1}^I \mu_{i} N_{i,t} = Q_t + \mathcal{H}_t + A_t h_t^\alpha - \xi_t \frac{h_t^{1+\nu}}{1+\nu},
\end{equation}
where $\mathcal{H}_t = \sum_{i=1}^I \mu_i \mathcal{H}_{i,t}$ and using the market clearing condition for stocks and goods.

Using the pricing condition for stocks and bonds, we obtain
\begin{align}
    \sum_{i=1}^I \mu_{i} N_{i,t} &= \mathbb{E}_t\left[ \sum_{k=1}^\infty \Lambda_{t,t+k} \pi_{t+k}\right] +\mathbb{E}_t\left[ \sum_{k=1}^\infty \Lambda_{t,t+k} \left( W_{t+k} h_{t+k} - \xi_{t+k} \frac{h_{t+k}^{1+\nu}}{1+\nu} \right)\right]+ A_t h_t^\alpha - \xi_t \frac{h_t^{1+\nu}}{1+\nu},\nonumber\\
    &= \mathbb{E}_t\left[ \sum_{k=0}^\infty \Lambda_{t,t+k} \left( A_{t+k} h_{t+k}^\alpha -\xi_{t+k} \frac{h_{t+k}^{1+\nu}}{1+\nu}  \right)\right] = A_{t-1} P_{t}.
\end{align}

We can write the market clearing for goods as follows:
\begin{align}
    \sum_{i=1}^I \frac{\mu_i N_{i,t}}{A_{t-1} P_t} \frac{C_{i,t}}{N_{i,t}} = \frac{A_t h_t^\alpha - \xi_t \frac{h_t^{1+\nu}}{1+\nu}}{A_{t-1} P_t}.
\end{align}

Using the definition of $\eta_{i,t}$  and the result $\sum_{i=1}^I \mu_{i} N_{i,t}  = A_{t-1} P_t$, we obtain the market clearing condition for goods in recursive notation:
\begin{equation}
    \sum_{i=1}^I \eta_{i} c_{i}(X,s) = \frac{x_s h(E)^\alpha - \xi \frac{h(E)^{1+\nu}}{1+\nu}}{P(X,s)}.
\end{equation}

Multiplying the market clearing condition for stocks by $Q_t$ and using the expression for $Q_t S_{i,t}$ given in Lemma \ref{lemma: investor_problem}, we obtain
\begin{equation}
    \sum_{i=1} \mu_i \omega_{i,t} (N_{i,t} - \tilde{C}_{i,t})  = \omega_{e,t} Q_t + \omega_{h,t}\mathcal{H}_{t},
\end{equation}
using the fact that $\mathcal{H}_{i,t} = \mathcal{H}_t$ and defining $\omega_{h,t} \equiv \omega_{h_i,t}$.

We will show next that we can write the right-hand side in the expression above as $ \omega_{e,t} Q_{t} + \omega_{h,t} \mathcal{H}_t =  Q_t +\mathcal{H}_t$. 
To show this fact, we start by the considering the return of holding the aggregate amount of stocks and human wealth:\small
\begin{align}
    R_{e,t+1} Q_{t} + R_{h,t+1} \mathcal{H}_t &= \left[Q_{t+1} + A_{t+1} h_{t+1}^\alpha  - W_{t+1} h_{t+1}\right] + \left[ \mathcal{H}_{t+1} + W_{t+1} h_{t+1} - \xi_{t+1} \frac{h_{t+1}^{1+\nu}}{1+\nu}\right] \nonumber\\
    &= A_t P_{t+1} = R_{r,t+1} \left[ A_{t-1} P_{t} - \left( A_t h_t^\alpha - \xi_t \frac{h_t^{1+\nu}}{1+\nu}  \right) \right]\nonumber\\
    &= R_{r,t+1} \left( Q_t + \mathcal{H}_t \right).
\end{align}\normalsize

This implies the following condition among excess returns:
\begin{align}
    R_{e,t+1}^e Q_{t} + R_{h,t+1}^e \mathcal{H}_t &= R_{r,t+1}^e \left( Q_t + \mathcal{H}_t \right)\nonumber\\
    R_{r,t+1}^e \omega_{e,t} Q_{t} + R_{r,t+1}^e \omega_{h,t} \mathcal{H}_t &= R_{r,t+1}^e \left( Q_t + \mathcal{H}_t \right),
\end{align}
which gives the desired result $ \omega_{e,t} Q_{t} + \omega_{h,t} \mathcal{H}_t =  Q_t +\mathcal{H}_t$, where we used the fact that $R_{k,t+1}^e = \omega_{k,t} R_{r,t+1}^e$ for $k \in \{e,h\}$. 

Using the result above, we can write the market clearing condition on stocks as follows:
\begin{align}
    \sum_{i=1} \frac{\mu_i N_{i,t} (1- c_{i,t} )}{Q_t + \mathcal{H}_t} \omega_{i,t}  &= 1\nonumber\\
    \sum_{i=1} \frac{\eta_{i,t} A_{t-1}P_t (1- c_{i,t} )}{A_{t-1} P_t - \left(A_t h_{t}^\alpha - \xi_t \frac{h_t^{1+\nu}}{1+\nu} \right)} \omega_{i,t}  &= 1.
\end{align}
    
Using the fact that $\frac{A_t h_{t}^\alpha - \xi_t \frac{h_t^{1+\nu}}{1+\nu}}{A_{t-1} P_t} = \sum_{i=1}^I \eta_{i,t} c_{i,t}$, we obtain
\begin{align}
    \sum_{i=1} \tilde{\eta}_{i,t} \omega_{i,t}  &= 1,
\end{align}
where $\tilde{\eta}_{i,t} \equiv \frac{\eta_i(1-c_{i,t})}{\sum_{j=1}^I\eta_j(1-c_{j,t})}$.

